\documentclass[10pt,a4paper]{article}
\usepackage[margin=2.1cm]{geometry}
\usepackage{booktabs}
\usepackage{amsmath}
\usepackage{microtype}
\usepackage[hidelinks]{hyperref}
\usepackage{enumitem}
\usepackage{xcolor}

\setlength{\parskip}{3pt}
\setlist{nosep,leftmargin=1.4em}
\newcommand{\ans}[1]{\par\smallskip\noindent\hspace*{1.2em}\begin{minipage}{0.94\linewidth}\small\itshape #1\end{minipage}\par\smallskip}

\title{\vspace{-1.4cm}\textbf{RAG Assignment: Baseline vs.\ Re-Ranking}}
\author{\texttt{<Name>} \quad$\cdot$\quad \texttt{<Roll number>}}
\date{October 2026}

\begin{document}
\maketitle
\vspace{-0.9cm}

\section*{Setup}

\textbf{Document.} Apple Inc.'s Form 10-K for fiscal 2024: 121 pages, about 64,000 words.
I chose it because it mixes very different kinds of text. Pages 1--60 are the annual report
proper --- financial statements, segment results, risk factors. Pages 61--121 are filed
exhibits: debt indentures, employee stock plans, an insider trading policy and a
compensation recovery (clawback) policy. Words like \emph{hedge}, \emph{pledge} and
\emph{compensation} mean different things in different parts of the document, which is
exactly what a bi-encoder struggles with.

\textbf{Chunking.} 547 chunks of at most 200 tokens with a 40-token overlap, packed sentence
by sentence. I sized chunks in tokens rather than words after finding that 150-word chunks
left 6.9\% of them longer than the bi-encoder's 256-token limit. Those chunks would have
been silently truncated for the bi-encoder but not for the cross-encoder (limit 512),
biasing the comparison in the re-ranker's favour. After the change, no chunk exceeds 229
tokens.

\begin{center}\small
\begin{tabular}{lll}
\toprule
 & \textbf{Pipeline 1 --- Baseline} & \textbf{Pipeline 2 --- Re-ranking} \\
\midrule
Stage 1 & \texttt{all-MiniLM-L6-v2}, Top 10 by cosine (FAISS) & same Top 10 \\
Stage 2 & --- & \texttt{ms-marco-MiniLM-L-6-v2} scores all 10 pairs \\
Context & Top 3 by cosine & Top 3 by cross-encoder score \\
LLM & \texttt{qwen2.5:3b} (Ollama), temperature 0 & identical model and prompt \\
\bottomrule
\end{tabular}
\end{center}

The prompt tells the model to answer only from the context and to say so if the answer is
missing. Since the candidate pool, prompt and model are identical, any difference between
Answer 1 and Answer 2 comes from which three chunks were chosen. Everything ran on CPU.

% ---------------------------------------------------------------------
\section{Architecture comparison}

A \textbf{bi-encoder} embeds the query and each chunk \emph{separately} into fixed vectors and
scores them by cosine similarity. Chunk vectors are computed once, offline, so a search costs
one query embedding plus a fast vector lookup. The weakness is that each chunk is compressed
into a single vector before the query is known. Everything the chunk says gets averaged
together, and the model cannot weigh one word of the chunk against one word of the query.
\emph{``Employees may not hedge Apple stock''} and \emph{``Apple hedges foreign currency
exposure''} end up close in vector space because both are about hedging.

A \textbf{cross-encoder} concatenates the query and the chunk into one input and runs them
through the transformer \emph{together}, so self-attention connects every query token with
every chunk token. It can notice that the query asks about employees and the chunk is about
the company, or that a ``not'' changes the meaning. It outputs a relevance score rather than
an embedding, so nothing can be precomputed: it must run once per (query, chunk) pair at
query time. That makes it more accurate and much more expensive, which is why it only
re-ranks the bi-encoder's ten candidates.

% ---------------------------------------------------------------------
\section{Qualitative and rank-shift analysis}

Each query has a gold chunk --- the one that actually contains the answer, found by reading
the document. Table~\ref{tab:ranks} shows where it ranked in each pipeline.

\begin{table}[h]
\centering\small
\caption{Rank of the gold chunk. Only the Top 3 reach the LLM.}\label{tab:ranks}
\begin{tabular}{llcc l}
\toprule
 & Query (category) & Baseline & Re-ranked & Outcome \\
\midrule
Q1 & Who serves as the Insider Trading Compliance Officer? \textit{(buried)} & \#7 & \#1 & rescued \\
Q2 & Is iPhone's share of net sales the same in every region? \textit{(exception)} & \#5 & \#1 & rescued \\
Q3 & Are puts and calls on Apple stock restricted for insiders? \textit{(keyword)} & \#7 & \#1 & rescued \\
Q4 & Which employees are not allowed to pledge Apple stock? \textit{(negation)} & \#7 & \#2 & rescued \\
Q5 & How is incentive-based compensation defined for clawback? \textit{(keyword)} & \#5 & \#1 & rescued \\
\midrule
Q6 & Can executive officers pledge Apple stock as collateral? \textit{(control)} & \#1 & \#1 & both fine \\
Q7 & Why did sales in China fall in 2024? \textit{(control)} & \#1 & \#1 & both fine \\
\midrule
Q8 & How many full-time equivalent employees did Apple have? & $>$10 & $>$10 & neither \\
Q9 & Can shares be pledged? & $>$10 & $>$10 & neither \\
Q10 & When is Apple excused from clawing back executive pay? & \#7 & \#9 & neither \\
\bottomrule
\end{tabular}
\end{table}

\subsection*{Where re-ranking fixed the answer (Q1--Q5)}

In all five rescue queries the gold chunk sat at rank 5--7, inside the Top 10 but outside
the Top 3, so the baseline never showed it to the LLM. Answer 1 was correct in \textbf{0 of 5}
cases and Answer 2 in \textbf{4 of 5}.

\textbf{Omitted answer (Q1, Q3, Q4).} The baseline's context did not contain the answer, and
the model correctly refused. For Q4 the baseline's Top 3 were a line from the exhibit index,
the closing lines of the clawback policy and a description of the employee stock plan ---
chunks that match ``employees'' and ``executive officers'' but say nothing about pledging.
The re-ranker promoted the insider trading policy from \#7 to \#2:
\ans{Answer 1: ``The context does not contain this information.''\\[2pt]
Answer 2: ``Members of Apple's Board of Directors and executive officers are prohibited from
pledging Apple stock as collateral.''}
Q4 is also the negation case. The question asks who is \emph{not} allowed to do something;
the cross-encoder read the query and the prohibition together and ranked the matching rule
second, which the bi-encoder had placed seventh. Q3 is a distractor case of a different kind. The baseline's Top 3 were
other sections of the same insider trading policy, on material non-public information and
on trading plans. They share the query's vocabulary (``insiders'', ``Apple stock'',
``restricted'') but none mentions puts, calls or derivatives. ``Puts'' appears in exactly one
chunk of the whole document, in a list of prohibited instruments, and the cross-encoder found
it.

\textbf{Misstated rule (Q5).} This was the most useful case, because the baseline answer
\emph{looked} right. It is also a keyword trap: the baseline's \#2 and \#3 were ``Recoupment''
clauses from employee stock award agreements, which overlap heavily with ``clawback'' and
``compensation'' but are not the clawback policy at all. With the gold chunk at \#5, the
baseline answered from its \#1, a neighbouring chunk of the policy that describes which
awards are in scope and when compensation counts as ``received'':
\ans{Answer 1: ``Incentive-Based Compensation is defined as awards granted to a Covered
Executive that are outstanding and unvested as of October 2, 2023, and is `Received' in the
fiscal period that the relevant financial reporting measure is attained\ldots''\\[2pt]
Answer 2: ``Incentive-Based Compensation for clawback purposes includes any compensation
granted, earned, or vested based in whole or in part upon the attainment of one or more
financial reporting measures\ldots''}
Answer 1 is fluent and uses the right vocabulary, but it presents a scope rule as the
definition. An automated check for the phrase ``financial reporting measure'' marked it
correct, so every answer in the notebook was also labelled by hand.

\subsection*{Where re-ranking did not help (Q2, Q8--Q10)}

\textbf{Retrieval fixed, generation still failed (Q2).} The re-ranker moved the right chunk
to \#1, but both answers said the context did not contain the information. The chunk states
that the revenue mix was ``generally consistent for each reportable segment\ldots except in
Greater China, where iPhone revenue represented a moderately higher proportion,'' wedged
between sentences about deferred revenue. The 3B model did not connect that to the question.
Re-ranking is necessary here, but not sufficient.

\textbf{Recall ceiling (Q8).} The gold chunk is not in the Top 10 even when the query uses
the document's exact phrase, ``full-time equivalent employees.'' The figure (164,000) is one
sentence inside a chunk mostly about compensation and benefits, so the chunk's averaged
vector is about compensation. The cross-encoder can only reorder the ten candidates it is
given; it cannot recover a chunk the bi-encoder never retrieved.

\textbf{Re-ranking made the answer worse (Q9).} The document contains four near-identical
restricted stock unit agreements, each saying unvested shares ``shall not be sold, assigned,
transferred, \emph{pledged}\ldots''. These crowded the real answer out of the Top 10, and the
cross-encoder promoted two of them into the Top 3 because they match ``pledged'' so closely.
The baseline had refused; the re-ranked answer contradicts itself:
\ans{Answer 2: ``Yes, shares can be pledged as per the restrictions on transfer mentioned in
the context. Shares of Common Stock are not allowed to be sold, assigned, transferred,
pledged\ldots''}
This is negation failing in the LLM, not the retriever: it saw ``pledged'' and answered yes.

\textbf{Hallucination in both pipelines (Q10).} The gold chunk says recovery can be excused
if it would be ``excessively burdensome,'' would cost more than it recovers, or would violate
foreign law. It ranked \#7 in the baseline and fell to \#9 after re-ranking. Without it, both
answers stated the opposite with confidence:
\ans{Answer 1: ``Apple is not excused from clawing back executive pay under this
policy\ldots The policy does not provide a specific exemption or exception for when Apple is
excused from such actions.''}
This is the worst outcome in the set: a fluent, specific answer that is backwards. It also
shows that a confident ``no exception exists'' from a RAG system can simply mean the
exception was never retrieved.

% ---------------------------------------------------------------------
\section{Latency profiling}

Each stage was timed over 20 repeats for each of the 10 queries (200 runs) after a warm-up
pass. The LLM was timed separately on the five rescue queries. Ollama caches processed
prompts, including from earlier runs, so repeated prompts looked about ten times faster than
they really are; each timing call was given a unique prefix so that it had to be processed
from scratch.

\begin{center}\small
\begin{tabular}{lrr}
\toprule
Stage (median, ms) & Pipeline 1 & Pipeline 2 \\
\midrule
Query embedding (bi-encoder) & 12.1 & 12.1 \\
FAISS search over 547 vectors & 0.4 & 0.4 \\
Cross-encoder, 10 pairs & --- & 397.4 \\
\midrule
\textbf{Retrieval total} & \textbf{12.5} & \textbf{409.0} \\
\midrule
LLM: reading the prompt & 20,864 & 22,107 \\
LLM: generating the answer & 727 & 3,493 \\
LLM total, incl.\ request overhead & 21,642 & 29,149 \\
\midrule
\textbf{End to end (retrieval + LLM)} & \textbf{21,655} & \textbf{29,558} \\
\bottomrule
\end{tabular}
\end{center}

Re-ranking makes retrieval about \textbf{33 times slower} (12.5 ms to 409 ms), but that is
only \textbf{1.3\%} of the end-to-end time, because a CPU-only LLM spends about 21 seconds
just reading a three-chunk prompt. The larger end-to-end gap comes from answer length, not
retrieval: Pipeline 2 produced real answers (median 38 tokens) where Pipeline 1 mostly said
the context lacked the information (9 tokens). With a GPU-served or API model, generation
would take a second or two and the 400 ms re-ranking cost would become a much larger share.

% ---------------------------------------------------------------------
\section{Scaling: why not cross-encode the whole collection?}

Because the cross-encoder's cost grows with the number of chunks, and none of it can be done
ahead of time. Every chunk must be read together with the query at query time. The
bi-encoder does the expensive work once, offline; at query time it embeds one short query
and compares vectors, which is a dot product per chunk.

I measured both on this document. Cross-encoding all 547 chunks for one query took
\textbf{22.7 seconds} (41.5 ms per pair), against 0.36 s for the ten candidates. One
cross-encoder pair costs as much as about 63,000 vector comparisons.

\begin{center}\small
\begin{tabular}{rrr}
\toprule
Chunks in collection & Cross-encode everything & Two-stage (flat index) \\
\midrule
547 & 22.7 s & 0.37 s \\
10,000 & 6.9 min & 0.38 s \\
100,000 & 1.2 h & 0.44 s \\
1,000,000 & 11.5 h & 1.03 s \\
\bottomrule
\end{tabular}
\end{center}

These are linear extrapolations from the measured per-pair and per-vector costs. The
two-stage column uses an exact flat index, which still grows linearly but very slowly;
approximate indexes such as HNSW grow sub-linearly. Even for this single document, a
22-second wait per query is too slow for interactive use, and at a million chunks it would
take half a day. The two-stage design uses the cheap model to narrow the search and the
expensive model only where precision matters.

% ---------------------------------------------------------------------
\section*{Takeaways}
\begin{itemize}
\item Re-ranking helped most when the right chunk was already in the Top 10 but buried by
chunks that shared its vocabulary. It took the right answers from 0/5 to 4/5 on those queries.
\item It cannot fix recall. If the bi-encoder misses the chunk (Q8, Q9), the cross-encoder
never sees it, so the retrieval depth of the first stage matters as much as the re-ranker.
\item Lexical decoys can fool the re-ranker too (Q9), and a missing exception produced the
most dangerous answer in the set (Q10).
\item On CPU the re-ranking cost (about 0.4 s) is small next to LLM generation, so the
accuracy gain is clearly worth it here.
\end{itemize}

\vspace{0.2cm}
\noindent\small Code, outputs and data: \texttt{rag\_baseline\_vs\_rerank.ipynb} in the
repository. Temperature-0 generation was not word-for-word identical across three full runs
(14 of 20 answers identical), but every answer's label held in all three.

\end{document}
